\section{Bernoulli graphings and finite permutation models}
\label{models:section}

We next turn the relative graphing quantity into a generator bound for
explicitly presented groups. The finite models used here are actions of
$A$ alone. We then construct models whose defining rows have the
expansion and overlap properties needed for the rank argument.

A finite exact right $A$-set is a finite set equipped with a right
action of $A$ by permutations; the action laws hold as equalities.
For such a set $V$, put
\begin{equation}
 D_V=\left\langle x_v\ (v\in V)\ \middle|\
 x_vx_{vu_1}\cdots x_{vu_{99}}=1\quad(v\in V)\right\rangle.
 \label{models:presentation}
\end{equation}
Here $u_0=1$ and $u_i=u_{i-1}ab_i$, as in the definition of $J$.
The notation $\operatorname{rk}(D)$ means the minimum size of a
generating set of a group $D$. The groups $D_V$ need not be finite.
We write $\operatorname{Fix}_V(g)=\{v\in V:vg=v\}$.

\subsection{The transfer inequality}

\begin{proposition}\label{models:transfer}
Let $V_k$ be finite exact right $A$-sets, with $n_k=|V_k|\to\infty$,
such that for every fixed $g\in A\setminus\{1\}$,
\begin{equation}
 \frac{|\{v\in V_k:vg=v\}|}{n_k}\longrightarrow0.
 \label{models:fixed-points}
\end{equation}
For the free Bernoulli action on $X$ from
Section~\ref{actions:section}, whose cost agrees with the full product
action on $X_0=[0,1]^\Gamma$, one has
\[
 \delta_X\ge\limsup_{k\to\infty}
 \frac{\operatorname{rk}(D_{V_k})}{n_k}.
\]
\end{proposition}

\begin{proof}
We use right actions, so that $(xg)(h)=x(gh)$. Choose representatives
$h_c$ for the cosets $Ah_c$ in $\Gamma$, and let
$\Omega=[0,1]^{A\backslash\Gamma}$ with its product probability measure.
The coordinate rearrangement
\[
 x\longmapsto \bigl((x(gh_c))_c\bigr)_{g\in A}
\]
identifies the restriction of $X_0$ to $A$ with $\Omega^A$, with independent
base coordinates and right action $(\xi q)(g)=\xi(qg)$.
We write $\mu$ for the product measure on this space. All graphing
identities are considered on invariant conull subsets, so the
restriction to $X$ does not change the argument.

For every finite exact $A$-set $V$, define
$\theta:V\times A\to D_V$ on the free basis of $A$ by
\[
 \theta(v,a)=x_v,\qquad \theta(v,b_i)=1,
 \qquad
 \theta(v,s^{-1})=\theta(vs^{-1},s)^{-1}
\]
for each basis letter $s$. Extend by multiplying these values along a
word. Cancellation of adjacent inverse letters cancels their two
values at the corresponding sources. Since $A$ is free, this gives a
well-defined function satisfying
\begin{equation}
 \theta(v,gh)=\theta(v,g)\theta(vg,h).
 \label{models:cocycle}
\end{equation}
The word $w=u_{99}a=ab_1ab_2\cdots ab_{99}a$ gives
\[
 \theta(v,w)=x_vx_{vu_1}\cdots x_{vu_{99}}=1.
\]
Also $\theta(v,b_i)=1$ at every source. The inverse rule and
\eqref{models:cocycle} therefore give
\begin{equation}
 \theta(v,j)=1\qquad(v\in V,\ j\in J).
 \label{models:j-trivial}
\end{equation}
Notice that $\theta(v,g)$ depends on the group label $g$, not merely on
the pair of endpoints $(v,vg)$.

Fix a finite-cost graphing $\mathcal E$ such that
$\mathcal R_J\vee\mathcal E=\mathcal R_A$, and fix $\epsilon>0$.
Such finite-cost graphings suffice for the infimum defining $\delta_X$:
the full $a$-translation, together with $\mathcal R_J$, already
generates $\mathcal R_A$.
On an invariant conull free set, split each map in $\mathcal E$
according to its unique displacement in $A$. This produces countably
many restrictions
\[
 \varphi_\ell:L_\ell\longrightarrow L_\ell g_\ell,
 \qquad \varphi_\ell(x)=xg_\ell,
 \qquad g_\ell\in A,
\]
with $\sum_\ell\mu(L_\ell)=C(\mathcal E)$.

A path type is a finite list of full $J$-translations and signed
indices $\ell$ of these paid maps. Its label is the product in $A$ of
its successive step labels. Almost every $x$ has a valid path to $xa$.
Freeness implies that every such path has label exactly $a$ in $A$.
There are countably many path types. Choose finitely many, each with
label $a$, whose validity sets cover a set of measure at least
$1-\epsilon$. Let $I$ be their finite set of paid indices, and let
$m_\ell$ count all occurrences of index $\ell$ among these types.

Choose cylinder-measurable sets $L'_\ell$, depending on finitely many
$A$-coordinates, such that
\[
 \sum_{\ell\in I}(1+m_\ell)
       \mu(L_\ell\mathbin{\triangle}L'_\ell)<\epsilon.
\]
This approximation is available because finite-coordinate measurable
sets form an algebra generating the product sigma-algebra. Its closure
for the metric $\mu(E\mathbin{\triangle}F)$ is a sigma-algebra, so
every measurable set is approximable in this metric.

For clarity, a forward step of index $\ell$ after a prefix with label
$h$ tests $xh\in L_\ell$. An inverse step after that prefix tests
$xhg_\ell^{-1}\in L_\ell$. Replace these tests by the corresponding
tests in $L'_\ell$, keeping all path labels unchanged. Each occurrence
changes its validity predicate on a set of measure at most
$\mu(L_\ell\mathbin{\triangle}L'_\ell)$, by invariance of $\mu$.
If $P'$ is the union of the modified path validity sets, then
\begin{equation}
 \sum_{\ell\in I}\mu(L'_\ell)\le C(\mathcal E)+\epsilon,
 \qquad \mu(P')\ge1-2\epsilon.
 \label{models:cylinder-bounds}
\end{equation}
All sets $L'_\ell$ and $P'$ depend on a single finite coordinate set
$F\subset A$: for example, a test of an $F_\ell$-cylinder after prefix
$h$ uses coordinates $hF_\ell$, or $hg_\ell^{-1}F_\ell$ for an inverse
step. This finite set is fixed before $k$ tends to infinity.

We have reduced generation to finitely many cylinder tests. We now
evaluate these same tests on $V_k$.
For independent base labels $(z_v)_{v\in V_k}$ with law $\Omega$, put
\[
 \xi_v=(z_{vg})_{g\in A}.
\]
Exactness of the action gives $\xi_{vq}=\xi_vq$. If $vg\ne vh$ for
all distinct $g,h\in F$, the $F$-coordinates of $\xi_v$ have exactly
the product distribution. The fraction of exceptional sources is at
most
\[
 \rho_k(F):=
 \sum_{\substack{g,h\in F\\g\ne h}}
 \frac{|\{v\in V_k:vgh^{-1}=v\}|}{n_k},
\]
which tends to zero by \eqref{models:fixed-points}. Consequently, for
every measurable $F$-cylinder $E$,
\begin{equation}
 \left|\frac1{n_k}\sum_{v\in V_k}
       \Pr(\xi_v\in E)-\mu(E)\right|\le\rho_k(F).
 \label{models:local-law}
\end{equation}
This uses no independence between different sources $v$.

For any realization of the labels, make a list containing
$\theta(v,g_\ell)$ for every $\ell\in I$ and every $v$ with
$\xi_v\in L'_\ell$. If $\xi_v\in P'$, follow a modified valid path
from $v$. A forward paid factor occurs on this list. For an inverse
paid step after prefix $h$, the factor is
\[
 \theta(vh,g_\ell^{-1})
   =\theta(vhg_\ell^{-1},g_\ell)^{-1},
\]
and its inverse-domain test puts
$\theta(vhg_\ell^{-1},g_\ell)$ on the list.
Each $J$-step contributes $1$ by \eqref{models:j-trivial}. Since the
path label is exactly $a$, \eqref{models:cocycle} expresses
$x_v=\theta(v,a)$ as a product from this list and its inverses.
The finite configuration $\xi_v$ need not be free; freeness was used
only to select the original path labels.

Add $x_v$ for every source with $\xi_v\notin P'$. The resulting list
generates $D_{V_k}$, so for every labeling,
\[
 \operatorname{rk}(D_{V_k})\le
 \sum_{\ell\in I}\sum_{v\in V_k}\mathbf1_{L'_\ell}(\xi_v)
 +\sum_{v\in V_k}\mathbf1_{(P')^c}(\xi_v).
\]
Taking expectations, using \eqref{models:local-law}, and then letting
$k\to\infty$ gives
\[
 \limsup_k\frac{\operatorname{rk}(D_{V_k})}{n_k}
 \le\sum_{\ell\in I}\mu(L'_\ell)+1-\mu(P')
 \le C(\mathcal E)+3\epsilon.
\]
Let $\epsilon$ tend to zero and then take the infimum over
$\mathcal E$. All finite path and cylinder tests were fixed before
taking the limit in $k$.
\end{proof}

\subsection{Expanding rows with few overlaps}

Regard the defining relator for $v$ as an ordered row with columns
$vu_0,vu_1,\ldots,vu_{99}$. For $T\subseteq V$, write
\[
 N(T)=\bigcup_{v\in T}\{vu_0,vu_1,\ldots,vu_{99}\}.
\]
A row is called bad if it has a repeated column or shares at least
two distinct columns with a different row. Let $b(V)$ be the number
of bad rows.

\begin{proposition}\label{models:sequence}
Put $K=e^{96}96^{99}/95^{95}$. Fix
\[
 0<\alpha<1/200,\qquad K\alpha^3<1/2.
\]
There exists a sequence of finite exact right $A$-sets $V_k$, with
$n_k=|V_k|\to\infty$, satisfying \eqref{models:fixed-points} and
\begin{align}
 |N(T)|&>96|T|
 &&(T\subseteq V_k,\ 1\le |T|\le\alpha n_k),
 \label{models:expansion}\\
 b(V_k)&=o(n_k).
 \label{models:bad-rows}
\end{align}
\end{proposition}

\begin{proof}
First, $a,u_1,\ldots,u_{99}$ is a free basis of $A$. Indeed, the
substitutions defining $u_i$ have the inverse substitutions
\[
 b_i=a^{-1}u_{i-1}^{-1}u_i.
\]
Between abstract free groups on the two proposed bases, these formulas
give inverse homomorphisms: they recover $b_i$ directly and recover
$u_i$ by induction. Assign independent uniform permutations of
$\{1,\ldots,n\}$ to the 100 elements of this basis. This defines an
exact right $A$-action.

Fix an integer $s$ with $1\le s\le\lfloor\alpha n\rfloor$. If a set
$T$ of size $s$ violates \eqref{models:expansion}, then $N(T)$ is
contained in a set $C$ of size $96s$ containing $T$, since $u_0=1$.
Such a set exists because $96s<n$. For fixed $T$ there are
$\binom{n-s}{95s}$ possible sets $C$. For one uniform permutation
$\sigma$, the probability that $T\sigma\subseteq C$ is
\[
 \frac{(96s)(96s-1)\cdots(96s-s+1)}{n(n-1)\cdots(n-s+1)}
 \le \left(\frac{96s}{n}\right)^s.
\]
Using independence of the 99 permutations assigned to $u_1,\ldots,u_{99}$,
followed by the bounds $\binom nr\le(en/r)^r$, gives
\begin{align*}
 \Pr\bigl(\exists T:\ |T|=s,\ |N(T)|\le96s\bigr)
 &\le\binom ns\binom n{95s}
                 \left(\frac{96s}{n}\right)^{99s}\\
 &\le\left[K\left(\frac{s}{n}\right)^3\right]^s.
\end{align*}
The exponent is $99-1-95=3$. Set $q=K\alpha^3<1/2$.
For every fixed integer $R\ge1$, the sum of these bounds over all
allowed sizes is at most
\[
 \sum_{s=1}^{\min(R,\lfloor\alpha n\rfloor)}
       [K(s/n)^3]^s+\frac{q^{R+1}}{1-q}.
\]
First let $n\to\infty$ and then $R\to\infty$. Thus
\eqref{models:expansion} holds with probability tending to one.
This includes $s=1$; if $\lfloor\alpha n\rfloor=0$, the assertion is
vacuous.

We also need \eqref{models:fixed-points}. Fix a nonempty freely reduced
word of length $L$ in the random basis, and a starting vertex $v_0$.
Follow the word, exposing permutation entries only when needed. Until
the first repeated vertex, the current vertex has not occurred earlier
except as the endpoint of its incoming edge. A positive next letter
could request an already exposed domain entry only by immediately
inverting that incoming edge. A negative next letter could request an
already exposed range entry only in the same way. Both are excluded
by free reduction. At step $t$, therefore, the requested partner is
uniform among at least $n-t+1$ available vertices. At most $t$ of
these vertices have already been visited. For $n\ge L$,
\[
 \Pr(v_0\text{ is fixed by the word})
 \le\sum_{t=1}^L\frac{t}{n-t+1}
 \le\frac{L(L+1)}{2(n-L+1)}.
\]
The same bound holds for the expected fixed-point fraction. It tends
to zero for each fixed nonidentity element of $A$.

Here is an explicit simultaneous choice. Enumerate the nonidentity
elements of $A$ as $g_1,g_2,\ldots$. At stage $k$, take $n_k>n_{k-1}$
so large that expansion fails with probability at most $1/4$, and
\[
 \mathbb E\frac{|\operatorname{Fix}_{V}(g_j)|}{n_k}
       \le\frac1{4k^2}\qquad(1\le j\le k).
\]
Markov's inequality and a union bound show that the probability that
some $j\le k$ has fixed-point fraction greater than $1/k$ is at most
$1/4$. Choose an outcome with expansion and all these fixed-point
bounds. This gives \eqref{models:fixed-points} for every fixed element,
without requiring estimates uniform over words of increasing length.

Finally we deduce \eqref{models:bad-rows}. A repeated column in row
$v$ gives $vu_i=vu_j$ for some $i\ne j$, so $v$ is fixed by the
nonidentity word $u_i u_j^{-1}$. If two distinct columns are shared
between rows $v\ne y$, choose indices with
\[
 vu_i=yu_k,\qquad vu_j=yu_l.
\]
Then $i\ne j$ and $k\ne l$ by distinctness of the columns, while
$i\ne k$ and $j\ne l$ follow from injectivity of each permutation.
Consequently $v$ is fixed by
\[
 u_i u_k^{-1}u_lu_j^{-1}.
\]
This word is nontrivial. In the cyclic order $(i,k,l,j)$, neighboring
indices differ. If all indices are nonzero, the displayed word is
freely reduced. Removing one occurrence of $u_0=1$ joins letters of
the same sign and leaves a nonempty reduced word. Two zero indices
can occur only in opposite positions; the remaining two letters then
have the same sign. Three zero indices are impossible. Thus every
possible overlap gives a fixed point of a nonidentity word from a
finite list. The union of the fixed sets in this list, together with
the finite list for repeated columns, contains every bad row.
Equation~\eqref{models:fixed-points} now gives $b(V_k)=o(n_k)$.
\end{proof}

For later use, the sequence in Proposition~\ref{models:sequence}
eventually satisfies
\[
 \alpha n_k\ge1,\qquad 100b(V_k)\le\alpha n_k.
\]
These are the finite size and exceptional-row bounds used in the
deterministic rank argument. Proposition~\ref{models:transfer} applies
to this same sequence.
